\documentclass[twocolumn]{aastex631}
\begin{document}
\title{A residual reionization ionizing-photon-budget shortfall for star-forming galaxies at z\~{}9 (o32 systematic corner)}
\author{NebulaMind Lab (autonomous pipeline)}
\affiliation{NebulaMind Lab, \url{https://lab.nebulamind.net}}
\begin{abstract}
Reionization ionizing-photon-budget reconciliation at z\~{}9: star-forming galaxies require f\_esc=0.452 (+0.455/-0.231) to close the budget (Madau-Dickinson SFRD, log xi\_ion=25.5+/-0.15, clumping C=2-5, JWST-SFRD tail), versus indirect-proxy-inferred f\_esc=0.080 (+0.146/-0.051) from LzLCS O32/beta calibrations. Median delta(required-inferred)=+0.339 dex-frac (16-84\%: +0.085 to +0.793); 91\% of the systematic MC shows a shortfall. Conclusion: a genuine SHORTFALL remains, and the sign holds under both O32 and beta calibrations. The result is bounded by the xi\_ion x clumping x proxy-calibration systematic, not by statistics. Generated autonomously from public data (jwst) via the NebulaMind Lab runner: a bounded, reproducible, descriptive study.
\end{abstract}
\section{Introduction}
Recent studies have highlighted a potential crisis in the photon budget required for reionization, suggesting that current observations may not account for the necessary ionizing photons [Muñoz2024]. This discrepancy has sparked interest in reconciling the ionizing-photon-budget using literature-anchored calculations. Previous works have explored various aspects of this problem, such as the role of absorption-dominated reionization and its increased demands on ionizing sources [Davies2021], as well as the calibration of excursion set reionization models to conserve ionizing photons [Park2022]. However, a comprehensive analysis of the photon budget using established literature values is still needed.
\section{Data and method}
To address this issue, we employed a systematic approach based on published literature values. Specifically, we utilized the cosmic star formation rate density (SFRD) from Madau \& Dickinson's (2014) analytic fitting function and adopted the ionizing efficiency (xi\_ion) value of log xi\_ion=25.5±0.15. Additionally, we relied on the O32/beta f\_esc proxy calibrations from LzLCS studies [Chisholm+22, Flury+22; Simmonds+24]. Notably, our method did not incorporate any new observational or catalog data, focusing instead on reconciling existing literature values.
\section{Result}
Our analysis revealed that star-forming galaxies at z\~{}9 require an escape fraction (f\_esc) of 0.452 (+0.455/-0.231) to reconcile the reionization ionizing-photon-budget. This value is significantly higher than the indirect-proxy-inferred f\_esc=0.080 (+0.146/-0.051) derived from LzLCS O32/beta calibrations. The median difference between the required and inferred escape fractions was found to be +0.339 dex-frac, with 91\% of systematic Monte Carlo simulations indicating a shortfall in ionizing photons.
\begin{figure}\centering\includegraphics[width=\columnwidth]{result.png}\caption{A residual reionization ionizing-photon-budget shortfall for star-forming galaxies at z\~{}9 (o32 systematic corner)}\end{figure}
\section{Caveats}
It is essential to acknowledge the limitations of our approach. As an automated, single-selection, uncalibrated measurement, our result relies heavily on the accuracy of the adopted literature values and calibrations. Systematic uncertainties in these inputs may affect the outcome, and further research is needed to refine these estimates. Moreover, our analysis does not account for potential variations in galaxy properties or environmental factors that could influence the ionizing photon budget. Therefore, while our study provides valuable insights into the reionization photon budget crisis, it should be interpreted with caution and considered as a stepping stone for future investigations.
\section*{References}
\footnotesize
[Muoz2024] Muñoz et al. (2024). Reionization after JWST: a photon budget crisis?. bibcode 2024MNRAS.535L..37M \par [Davies2021] Davies et al. (2021). The Predicament of Absorption-dominated Reionization: Increased Demands on Ionizing Sources. bibcode 2021ApJ...918L..35D \par [Park2022] Park et al. (2022). Calibrating excursion set reionization models to approximately conserve ionizing photons. bibcode 2022MNRAS.517..192P \par [Duncan2015] Duncan et al. (2015). Powering reionization: assessing the galaxy ionizing photon budget at z < 10. bibcode 2015MNRAS.451.2030D \par [Madau2017] Madau (2017). Cosmic Reionization after Planck and before JWST: An Analytic Approach. bibcode 2017ApJ...851...50M

\end{document}
